\documentclass[11pt,a4paper]{article}
\usepackage[T1]{fontenc}
\usepackage[utf8]{inputenc}
\usepackage[italian]{babel}
\usepackage{lmodern,geometry,booktabs,tabularx,enumitem,xcolor,hyperref,listings}
\geometry{margin=23mm}
\definecolor{blu}{HTML}{153B62}
\hypersetup{colorlinks=true,linkcolor=blu,urlcolor=blu}
\setlength{\parindent}{0pt}
\setlength{\parskip}{6pt}
\title{\textbf{Algorithm Challenge --- Esempio metodologico}\\Compressione Run-Length Encoding (RLE)}
\author{IIS Europa Unita --- 4C INFO\\Prof. Christian Di Maria}
\date{}
\begin{document}
\maketitle
\noindent\textbf{Scopo:} mostrare come analizzare, progettare, implementare, collaudare e discutere un algoritmo in Java. RLE \emph{non} è tra i dodici algoritmi assegnati ai sei gruppi. Non va consegnato come terzo algoritmo.

\section{Il problema}
RLE sostituisce ogni sequenza consecutiva di caratteri uguali con una coppia \texttt{conteggio + lettera}. Il programma accetta stringhe da 0 a 100\,000 caratteri, soltanto lettere maiuscole \texttt{A--Z}. La restrizione evita ambiguità nella decodifica.

\begin{center}
\large\texttt{AAAABBBCCDAA} $\longrightarrow$ \texttt{4A3B2C1D2A}
\end{center}

\section{Esecuzione manuale}
Scorriamo la stringa da sinistra verso destra. Finché il carattere è uguale al precedente incrementiamo il conteggio; a ogni cambio emettiamo il gruppo appena terminato. L'ultimo gruppo va emesso dopo il ciclo.

\begin{tabularx}{\linewidth}{@{}lXl@{}}
\toprule
Gruppo & Decisione & Output progressivo\\
\midrule
\texttt{AAAA} & Cambio A/B: emetti 4A & \texttt{4A}\\
\texttt{BBB} & Cambio B/C: emetti 3B & \texttt{4A3B}\\
\texttt{CC} & Cambio C/D: emetti 2C & \texttt{4A3B2C}\\
\texttt{D} & Cambio D/A: emetti 1D & \texttt{4A3B2C1D}\\
\texttt{AA} & Fine stringa: emetti 2A & \texttt{4A3B2C1D2A}\\
\bottomrule
\end{tabularx}

\section{Perché la decodifica è corretta?}
Per ogni gruppo consecutivo, la compressione memorizza esattamente la quantità e il carattere. La decompressione ricrea le lettere nell'ordine. Verifichiamo la proprietà \texttt{decomprimi(comprimi(testo)).equals(testo)} sugli input ammessi. Non basta confrontare soltanto le lunghezze.

\section{Progettazione a oggetti}
\begin{tabularx}{\linewidth}{@{}lX@{}}
\toprule
Classe & Responsabilità\\
\midrule
\texttt{CompressioneRLE} & Convalida, compressione e decompressione.\\
\texttt{RisultatoCompressione} & Dati della prova, statistiche e \texttt{toString()}.\\
\texttt{Main} & Input tramite \texttt{Scanner} e visualizzazione.\\
\texttt{TestRLE} & 18 verifiche automatiche indipendenti.\\
\bottomrule
\end{tabularx}

L'esempio non forza un'ereditarietà artificiale: presenta una sola strategia. Nei progetti di gruppo si dovranno progettare due strategie e considerare un'eventuale generalizzazione significativa.

\newpage
\section{Testing}
Ogni risultato atteso deve poter essere ricostruito senza fidarsi del programma che si sta testando.

\begin{tabularx}{\linewidth}{@{}l l X@{}}
\toprule
Input & Output atteso & Obiettivo\\
\midrule
\texttt{AAAABBBCCDAA} & \texttt{4A3B2C1D2A} & Sequenze multiple\\
\texttt{AAAAAAAAAAAA} & \texttt{12A} & Conteggio a più cifre\\
\texttt{ABC} & \texttt{1A1B1C} & Espansione sfavorevole\\
\texttt{AABBAA} & \texttt{2A2B2A} & Gruppi separati\\
Stringa vuota & Stringa vuota & Caso limite\\
\texttt{A1} & Errore & Input non ammesso\\
\bottomrule
\end{tabularx}

Il programma contiene 18 test automatici che controllano output atteso e, per i casi validi, ricostruzione esatta. Consultare \texttt{tests/casi-di-test.md}.

\section{Confronto sperimentale}
\begin{tabularx}{\linewidth}{@{}X r r r@{}}
\toprule
Input & Prima & Dopo & Risparmio\\
\midrule
\texttt{AAAAAAAAAAAA} & 12 & 3 & 75,00\%\\
\texttt{AAAABBBCCDAA} & 12 & 10 & 16,67\%\\
\texttt{ABC} & 3 & 6 & -100,00\%\\
\bottomrule
\end{tabularx}

La misura riguarda i \emph{caratteri della rappresentazione}, non i byte occupati su disco. Un valore negativo indica che RLE ha prodotto una stringa più lunga. Spiegare il contesto della misura è parte del lavoro algoritmico.

\section{Esecuzione}
Dalla radice del repository con un JDK installato:
\begin{lstlisting}
javac -encoding UTF-8 -d out src/*.java
java -cp out Main
java -cp out TestRLE
\end{lstlisting}
Il test runner deve mostrare \texttt{Test superati: 18/18}. In IntelliJ IDEA aprire la cartella come progetto e configurare il JDK.

\section{Domande per riflettere}
\begin{enumerate}[leftmargin=*]
\item Perché bisogna emettere l'ultimo gruppo dopo il ciclo?
\item Che cosa cambierebbe se il testo potesse contenere cifre?
\item In quali casi RLE comprime e in quali espande?
\item Perché verificare la decodifica invece della sola lunghezza?
\item Come progetteresti nuovi test senza modificare l'algoritmo?
\end{enumerate}

\section{Indicazioni per i gruppi}
Ogni gruppo deve sviluppare i \textbf{due algoritmi assegnati} in un unico repository e consegnare su Moodle il link GitHub. Ciascuno dei tre studenti deve conoscere l'intero progetto, i test, le classi e le conclusioni. La valutazione prevede \textbf{40\% progetto comune e 60\% interrogazione individuale}.

\bigskip
\begin{center}\color{blu}\textbf{Conoscere un algoritmo significa saperlo spiegare, verificare e discutere.}\end{center}
\end{document}
